% Template Dissertação de Mestrado abnTeX2 — pesquisador-br-skill
% Compilação: pdflatex + bibtex + pdflatex + pdflatex
% Diferenças vs TCC: estrutura mais robusta, inclui errata, mais
% capítulos, preâmbulo de mestrado, banca de 3 examinadores.

\documentclass[
    12pt,
    openright,
    twoside,
    a4paper,
    chapter=TITLE,
    sumario=tradicional,
    section=TITLE,
    english,
    brazil
]{abntex2}

\usepackage{cmap}
\usepackage[utf8]{inputenc}
\usepackage[T1]{fontenc}
\usepackage{lmodern}
\usepackage{lastpage}
\usepackage{indentfirst}
\usepackage{color}
\usepackage{graphicx}
\usepackage{microtype}
\usepackage{lipsum}                 % gerador de texto fake
\usepackage[brazilian,hyperpageref]{backref}
\usepackage[alf]{abntex2cite}
\usepackage{amsmath, amssymb}       % matemática
\usepackage{booktabs}               % tabelas profissionais

\titulo{TÍTULO DA DISSERTAÇÃO: subtítulo se houver}
\autor{Nome Completo Sobrenome}
\local{Manaus -- AM}
\data{2026}
\orientador{Prof.\textsuperscript{a} Dr.\textsuperscript{a} Nome do Orientador}
% \coorientador{Prof. Dr. Nome do Coorientador}
\instituicao{Universidade Federal do Amazonas \par
            Programa de Pós-Graduação em [Nome do Programa] \par
            Mestrado em [Área]}
\tipotrabalho{Dissertação (Mestrado)}
\preambulo{Dissertação apresentada ao Programa de Pós-Graduação
em [Nome do Programa] da Universidade Federal do Amazonas, como
requisito parcial para a obtenção do título de Mestre em [Área].
\par
Linha de pesquisa: [Linha].}

\usepackage{hyperref}
\hypersetup{
    pdftitle={\@title},
    pdfauthor={\@author},
    colorlinks=true,
    linkcolor=blue,
    citecolor=blue,
    urlcolor=blue,
}

\begin{document}
\selectlanguage{brazil}
\frenchspacing

% PRÉ-TEXTUAIS
\pretextual
\imprimircapa
\imprimirfolhaderosto

% Folha de aprovação (3 examinadores: orientador + 2)
\begin{folhadeaprovacao}
\begin{center}
{\ABNTEXchapterfont\large\imprimirautor}
\vspace*{\fill}
\begin{center}
\ABNTEXchapterfont\bfseries\Large\imprimirtitulo
\end{center}
\vspace*{\fill}

\hspace{.45\textwidth}\begin{minipage}{.5\textwidth}
\imprimirpreambulo
\end{minipage}\vspace*{\fill}
\end{center}

Dissertação aprovada. Manaus -- AM, \today:

\assinatura{\textbf{\imprimirorientadorRotulo} \\ \imprimirorientador}
\assinatura{\textbf{Prof. Dr. Nome 1 -- Examinador Interno} \\ \\ Universidade Federal do Amazonas}
\assinatura{\textbf{Prof. Dr. Nome 2 -- Examinador Externo} \\ \\ [Outra Instituição]}
\end{folhadeaprovacao}

% Dedicatória
\begin{dedicatoria}
\vspace*{\fill}\centering
\textit{Dedicatória.}
\vspace*{\fill}
\end{dedicatoria}

% Agradecimentos (mestrado: agência financiadora obrigatória)
\begin{agradecimentos}
O presente trabalho foi realizado com apoio da Coordenação de
Aperfeiçoamento de Pessoal de Nível Superior -- Brasil (CAPES) --
Código de Financiamento 001 / processo nº [...].

Agradeço à orientadora, à banca, à família e aos colegas de
programa.
\end{agradecimentos}

% Epígrafe
\begin{epigrafe}
\vspace*{\fill}
\begin{flushright}
\textit{``Citação inspiradora.''} \\
(AUTOR, ANO, p. X)
\end{flushright}
\end{epigrafe}

% Resumo
\setlength{\absparsep}{18pt}
\begin{resumo}
\noindent A presente dissertação investigou [problema] no contexto
de [recorte], adotando abordagem [qualitativa/quantitativa/mista].
Como objetivos: (i) [obj1]; (ii) [obj2]; (iii) [obj3]. Foram
coletados [dados] mediante [instrumento], analisados conforme
[método]. Os achados indicaram [resultado principal], evidenciando
[contribuição]. As limitações incluem [limitação], indicando-se
[agenda futura].

\noindent\textbf{Palavras-chave}: termo 1; termo 2; termo 3;
termo 4; termo 5.
\end{resumo}

% Abstract
\begin{resumo}[Abstract]
\begin{otherlanguage*}{english}
\noindent This dissertation investigated [problem] in the context
of [scope], adopting a [qualitative/quantitative/mixed] approach.
The objectives were: (i) [obj1]; (ii) [obj2]; (iii) [obj3]. [Data]
was collected through [instrument] and analyzed using [method].
Findings indicated [main result], evidencing [contribution].
Limitations include [limitation], with [future agenda] suggested.

\noindent\textbf{Keywords}: term 1; term 2; term 3; term 4; term 5.
\end{otherlanguage*}
\end{resumo}

% Listas e sumário
\pdfbookmark[0]{\listfigurename}{lof}
\listoffigures*
\cleardoublepage

\pdfbookmark[0]{\listtablename}{lot}
\listoftables*
\cleardoublepage

% Lista de abreviaturas
\begin{siglas}
\item[CAPES] Coordenação de Aperfeiçoamento de Pessoal de Nível Superior
\item[CNPq] Conselho Nacional de Desenvolvimento Científico e Tecnológico
\item[CEP] Comitê de Ética em Pesquisa
\item[TCLE] Termo de Consentimento Livre e Esclarecido
\end{siglas}

\pdfbookmark[0]{\contentsname}{toc}
\tableofcontents*
\cleardoublepage

% TEXTUAIS
\textual

\chapter{INTRODUÇÃO}
\section{Contextualização}
\section{Problema de pesquisa}
\section{Objetivos}
\subsection{Geral}
\subsection{Específicos}
\section{Justificativa}
\section{Estrutura da dissertação}

\chapter{REFERENCIAL TEÓRICO}
\section{[Conceito 1]}
\section{[Conceito 2]}
\section{[Conceito 3]}
\section{Estado da arte}

\chapter{METODOLOGIA}
\section{Tipo e abordagem}
\section{Cenário/contexto}
\section{Sujeitos/participantes}
\section{Procedimentos de coleta}
\section{Procedimentos de análise}
\section{Aspectos éticos}

\chapter{RESULTADOS}
\section{[Eixo 1]}
\section{[Eixo 2]}
\section{[Eixo 3]}

\chapter{DISCUSSÃO}
\section{Implicações teóricas}
\section{Implicações práticas}
\section{Limitações}

\chapter{CONSIDERAÇÕES FINAIS}
\section{Síntese}
\section{Contribuições}
\section{Agenda futura}

% PÓS-TEXTUAIS
\postextual

\bibliography{referencias}

\begin{apendicesenv}
\partapendices
\chapter{ROTEIRO DE ENTREVISTA}
\chapter{TCLE -- TERMO DE CONSENTIMENTO}
\end{apendicesenv}

\begin{anexosenv}
\partanexos
\chapter{PARECER DO CEP}
\end{anexosenv}

\end{document}
